\documentclass[conference]{IEEEtran}
\usepackage{ctex}
\usepackage{cite}
\usepackage{amsmath,amssymb,amsfonts}
\usepackage{algorithmic}
\usepackage{graphicx}
\usepackage{textcomp}
\usepackage{xcolor}
\usepackage{booktabs}
\usepackage{multirow}
\usepackage{siunitx}
\def\BibTeX{{\rm B\kern-.05em{\sc i\kern-.025em b}\kern-.08em
    T\kern-.1667em\lower.7ex\hbox{E}\kern-.125emX}}

\begin{document}

\title{基于遗传算法的光电混合片上网络热感知任务映射}

\author{\IEEEauthorblockN{匿名作者}
\IEEEauthorblockA{双盲评审投稿}}

\maketitle

\begin{abstract}
微环谐振器（MRR）的热敏感性是光片上网络（ONoC）面临的核心挑战之一。局部温度变化会引起 MRR 谐振波长漂移，需要动态热调谐补偿；同时，当处理器核（PE）温度超过 DVFS 节流阈值时，计算时间被拉长，进一步影响系统性能和能耗。现有工作主要从器件级、电路级和路由级应对热管理问题，而\textbf{任务到 PE 的映射}作为设计时热源布局手段仍未被充分利用。本文提出一种基于遗传算法（GA）的热感知任务映射框架：每个候选映射通过\textbf{全系统 OMNeT++ 仿真}评估，该仿真包含 8$\lambda$ WDM 光层、SOA 放大、MRR 动态调谐、激光源、32 节点 RC 热网络和周期性 DVFS 反馈回路。GA 以基线归一化多目标代价函数为适应度，联合考虑峰值温度、温度标准差、热点数、makespan、通信距离、静态拥塞、DVFS 惩罚、负载均衡和总能耗。在覆盖 CCR 0.3--8 的四个基准上，所提方法均降低复合代价；其中 GEMM 和 MPEG4 同时改善热、性能、通信和能耗指标，VOPD 将 makespan 降低 \SI{50.5}{\percent}、总能耗降低 \SI{31.8}{\percent}，但牺牲部分热均匀性；HNN 则表现为明显的热-性能折中。结果表明，任务映射是 ONoC 热管理的有效设计时优化手段，但必须以多目标折中而非单一热指标最小化来建模。
\end{abstract}

\begin{IEEEkeywords}
光片上网络，热感知任务映射，遗传算法，微环谐振器，OMNeT++
\end{IEEEkeywords}

\section{引言}

光片上网络（ONoC）凭借硅光子波导和微环谐振器（MRR）提供的太比特级带宽密度、波分复用（WDM）能力和近乎与距离无关的数据移动能耗，已成为解决众核处理器电互连瓶颈的前景广阔方案~\cite{batten2008,bergman2018}。

然而，ONoC 的实际部署面临一个关键障碍：硅基 MRR 的强热光敏感性。对于典型 MRR，仅 \SI{1}{\degreeCelsius} 的温度偏差就会引起约 \SI{0.1}{\nano\meter} 的谐振波长漂移~\cite{padmaraju2014}，足以使 WDM 信道失谐并导致传输错误。在一个包含 2560 个 MRR 的 4$\times$4 Mesh ONoC 中，处理单元（PE）的局部发热——每个 PE 计算时功耗 \SI{2.5}{\watt}——会在芯片上产生超过 \SI{10}{\degreeCelsius} 的热梯度。为维持谐振对准，每个 MRR 需要逐器件热调谐，在最坏情况下消耗高达 \SI{5.12}{\watt} 的总功率。此外，当 PE 温度超过 DVFS 节流阈值时，计算时间被拉长，进一步恶化系统性能。

现有工作已从多个层面应对 ONoC 热挑战：器件级的无热化 MRR 设计~\cite{guha2012}、电路级的选择性调谐与热预算管理~\cite{zhou2013}、以及网络级的热感知路由~\cite{chittamuru2015}。然而，一个重要的设计时手段尚未得到充分利用：\textbf{任务到 PE 的映射}。计算任务在 PE 上的空间分布直接决定了芯片的热功率密度图。一个经过热优化的映射可以将发热任务分散到相距较远的 PE、减少通信引起的热耦合、并在任何运行时机制介入之前就最小化 DVFS 惩罚。

本文提出一种基于遗传算法（GA）的光电混合 NoC 热感知任务映射框架。主要贡献如下：

\begin{enumerate}
\item 一种 GA 优化方法，其中\textbf{每个个体通过全系统 OMNeT++ 仿真进行评估}，捕获完整的热-光耦合链：任务执行 $\rightarrow$ PE 发热 $\rightarrow$ DVFS 节流 $\rightarrow$ MRR 调谐 $\rightarrow$ 光路建立 $\rightarrow$ SOA 放大 $\rightarrow$ 性能与能耗；
\item 一个以静态基线为归一化参照的多目标代价函数，联合编码峰值温度、温度均匀性、热点数、makespan、通信距离、通信拥塞、DVFS 惩罚、负载均衡和总能耗，避免单纯追求降温而破坏性能；
\item 在覆盖计算通信比（CCR）0.3--8 的四个基准（GEMM、MPEG4、VOPD、HNN）上的实验验证，展示任务映射在不同工作负载下的收益、折中和优化边界。
\end{enumerate}

\section{系统架构}

\subsection{光电混合 NoC}

目标系统为 4$\times$4 Mesh ONoC，包含 16 个 TaskPE、16 个 5 端口光路由器和 1 个 GlobalBuffer（片外 DRAM 接口），如图~\ref{fig:arch} 所示。每个 PE 内含一个支持动态电压频率调节（DVFS）的任务执行核心。

控制面通信（任务分发、光路建立）经过电层 XY 维序路由的 Wormhole 交换网络。数据面通信使用光旁路：在 SETUP\_REQ/SETUP\_ACK 握手为路径分配波长后，数据微片通过 \texttt{sendDirect()} 经光波导直传，完全旁路电路由器。END\_FLIT 触发拆链和波长释放。

\begin{figure}[t]
\centering
% [待补充：架构示意图]
\caption{4$\times$4 光电混合 ONoC 架构：16 PE（灰色）、16 光路由器（蓝色）、GlobalBuffer（绿色）。实线箭头：电控制面。虚线箭头：光数据面旁路。}
\label{fig:arch}
\end{figure}

\subsection{光层}

每个光路由器采用 5$\times$5 微环交叉杆拓扑，每个交叉点串联 8 个 MRR（每 WDM 波长一个），每路由器共计 \textbf{160 个 MRR}，全芯片合计 \textbf{2560 个 MRR}。波长分配采用最低空闲优先策略，在 8 个 C 波段波长（1544.53--1555.75\,nm，ITU-T 200\,GHz 间隔）和 2 个空间信道上进行，每条边共 16 个波长槽。每条光电路占用 2 个波长，每边最多支持 8 条并发光电路，每波长数据率 256\,Gbps（PAM4，128\,GBaud）。

光功率由每个路由器输出端口后的半导体光放大器（SOA）提供，每个 SOA 消耗 \SI{80}{\milli\watt} 电泵浦功率。片外激光源提供每分支 \SI{1}{\milli\watt} 光发射功率（0\,dBm，经 16 路分光、\SI{3}{\decibel} 光栅耦合损耗和 \SI{1}{\decibel} 分光器额外损耗后），壁插效率 \SI{20}{\percent}，对应 \SI{5}{\milli\watt} 持续电功耗。

\subsection{热模型}

系统采用 32 节点双层 RC 热网络对芯片热行为建模~\cite{skadron2003}。16 个 PE 和 16 个路由器各为一个热节点，参数如下：
\begin{itemize}
\item \textbf{PE 节点}：$R_{\text{conv}}$ = \SI{8}{\kelvin\per\watt}（垂直对流），$R_{\text{lateral}}$ = \SI{10}{\kelvin\per\watt}（PE 间横向传导），$R_{\text{pe2router}}$ = \SI{3}{\kelvin\per\watt}（向下层路由器垂直耦合），$C_{\text{PE}}$ = \SI{1e-6}{\joule\per\kelvin}；
\item \textbf{路由器节点}：$R_{\text{conv}}$ = \SI{10}{\kelvin\per\watt}，$R_{\text{lateral}}$ = \SI{10}{\kelvin\per\watt}，$C_{\text{router}}$ = \SI{1e-7}{\joule\per\kelvin}。
\end{itemize}
环境温度 $T_{\text{amb}}$ = \SI{318.15}{\kelvin}（\SI{45}{\degreeCelsius}）。显式欧拉法求解器以 $\Delta t$ = \SI{100}{\nano\second} 步长推进。

每个时间步向热网络注入三类热源：
\begin{enumerate}
\item \textbf{PE 计算功耗}：活跃时 $P_{\text{PE}} = P_{\text{idle}} + P_{\text{compute}}$，经泄漏因子 $\exp((T - T_{\text{amb}})/15)$ 调制；
\item \textbf{SOA 泵浦功率}：每个 SOA \SI{80}{\milli\watt}，在光电路持续期间注入驱动路由器节点；
\item \textbf{MRR 调谐功率}：$P_{\text{tune},i} = \text{ringCount}_i \times \SI{0.5}{\milli\watt\per\nano\meter} \times \SI{0.10}{\nano\meter\per\kelvin} \times |T_i - T_{\text{amb}}|$，每个路由器根据自身温度和光路上活跃微环数量独立计算。
\end{enumerate}

\subsection{DVFS 反馈回路}

当 PE 瞬时温度超过 $T_{\text{throttle}}$ = \SI{327.15}{\kelvin}（\SI{54}{\degreeCelsius}）时，其计算时间被线性拉伸：$t_{\text{actual}} = t_{\text{nominal}} \times (1 + \beta \times (T - T_{\text{throttle}}))$，其中 $\beta = 0.1$（每 \si{\degreeCelsius} \SI{10}{\percent}）。此节流每 \SI{100}{\nano\second} 重新评估一次，形成温度与性能之间的闭环反馈。

\section{问题建模}

\subsection{任务图模型}

计算工作负载表示为有向无环任务图 $\mathcal{G} = (\mathcal{T}, \mathcal{E})$，每个节点 $\tau_i \in \mathcal{T}$ 包含属性：名义计算时间 $t_i^{\text{nom}}$、输出数据量 $d_i$ 和后继集 $\mathcal{S}_i$。边 $e_{ij} \in \mathcal{E}$ 表示从 $\tau_i$ 到 $\tau_j$ 的数据依赖。

NoC 拥有 $N = 16$ 个处理单元 $\{\text{PE}_0, \ldots, \text{PE}_{15}\}$，排列为 4$\times$4 网格。\textbf{任务映射}是一个函数 $M: \mathcal{T}_m \rightarrow \{0, \ldots, 15\}$，将每个可映射任务分配到一个 PE。GlobalBuffer 任务（数据注入/提取）不参与映射。

\subsection{代价函数}

给定完整映射 $M$ 及其 OMNeT++ 仿真输出，本文采用以静态基线 $M_0$ 为参照的归一化复合代价函数：

\begin{equation}
\begin{aligned}
\label{eq:cost}
\mathcal{F}(M) ={}& w_T f_T + w_\sigma f_\sigma + w_{\text{hot}} f_{\text{hot}}
+ w_M f_M + w_H f_H \\
&+ w_C f_C + w_D f_D + w_L f_L + w_E f_E .
\end{aligned}
\end{equation}

权重设为 $w_T=1.0$、$w_\sigma=1.0$、$w_{\text{hot}}=0.6$、$w_M=1.2$、$w_H=0.4$、$w_C=0.7$、$w_D=0.4$、$w_L=0.2$、$w_E=0.5$。这些权重不是为了将任一单项指标推到极限，而是使 GA 在热安全、执行时间、通信局部性和能耗之间搜索折中解。由于所有主要项都相对于同一工作负载的静态基线归一化，$\mathcal{F}(M)<\mathcal{F}(M_0)$ 表示相对基线的综合折中改善；不同目标函数版本之间的绝对代价值不可直接比较。

\subsubsection{热项}
\begin{equation}
f_T = \frac{\max(0, T_{\max}(M)-T_{\text{amb}})}
{\max(0, T_{\max}(M_0)-T_{\text{amb}})}
\end{equation}
$T_{\max}$ 是从 \texttt{pe-die-temperature} 矢量中提取的所有 PE、所有时间步真实峰值温度。为避免只优化单个最热点，代价函数同时加入温度均匀性项和热点数项：
\begin{align}
f_\sigma &= \frac{\sigma_T(M)}{\sigma_T(M_0)},\\
f_{\text{hot}} &=
\begin{cases}
N_{\text{hot}}(M)/N_{\text{hot}}(M_0), & N_{\text{hot}}(M_0)>0,\\
N_{\text{hot}}(M)/16, & N_{\text{hot}}(M_0)=0.
\end{cases}
\end{align}
其中 $\sigma_T$ 为全 PE、全时间步温度样本的标准差，$N_{\text{hot}}$ 为峰值温度超过 \SI{54}{\degreeCelsius} 的 PE 数。

\subsubsection{性能与 DVFS 项}
\begin{align}
f_M &= \frac{t_{\text{makespan}}(M)}{t_{\text{makespan}}(M_0)},\\
f_D &=
\begin{cases}
\eta_{\text{dvfs}}(M)/\eta_{\text{dvfs}}(M_0), & \eta_{\text{dvfs}}(M_0)>0,\\
\eta_{\text{dvfs}}(M)/100, & \eta_{\text{dvfs}}(M_0)=0.
\end{cases}
\end{align}
其中 $\eta_{\text{dvfs}}$ 是 16 个 PE 的平均节流比。显式加入 $f_M$ 是必要的：降低热点可能通过牺牲 DAG 并行度实现，若缺少 makespan 项，优化器会偏向过度冷却而产生不可接受的性能退化。

\subsubsection{通信与拥塞项}
\begin{equation}
f_H = \frac{\sum_{(i,j)\in\mathcal{E}} \text{hops}(M(\tau_i),M(\tau_j)) d_i}
{\sum_{(i,j)\in\mathcal{E}} \text{hops}(M_0(\tau_i),M_0(\tau_j)) d_i}
\end{equation}
其中 $\text{hops}(a,b)$ 为 4$\times$4 网格中的 Manhattan 距离。仅惩罚总跳数不足以刻画光层竞争，因此加入静态拥塞代理：
\begin{equation}
f_C = \frac{\max_{e\in\mathcal{P}} \text{bytes}_e(M)}
{\max_{e\in\mathcal{P}} \text{bytes}_e(M_0)} .
\end{equation}
$\text{bytes}_e$ 将任务图边按确定性 XY 路径投影到物理 mesh 边后累计字节量，近似反映波长槽争用风险。

\subsubsection{负载与能耗项}
\begin{equation}
f_L = \frac{\text{Var}(\text{load}_{\text{PE}_k}(M))/\text{ideal}^2}
{\text{Var}(\text{load}_{\text{PE}_k}(M_0))/\text{ideal}^2}, \qquad
f_E = \frac{E_{\text{total}}(M)}{E_{\text{total}}(M_0)} .
\end{equation}
其中 $\text{load}_{\text{PE}_k}$ 为分配到 PE$_k$ 的总名义计算时间，$E_{\text{total}}$ 包含 PE、SOA、MRR 调谐和激光器能耗。当某一基线分母为零时，实现中使用正值回退分母，避免通信为空或无 DVFS 的基准产生除零。

\section{遗传算法任务映射}

任务映射是一个在 $16^{|\mathcal{T}_m|}$ 个可能分配上的组合优化问题。本文提出的 GA 通过选择、交叉、变异和精英保留进化候选映射种群，以全系统 OMNeT++ 仿真作为适应度评估神谕。

\subsection{染色体编码与种群初始化}

每个个体编码为染色体 $\mathbf{c} = [p_1, p_2, \ldots, p_n]$，其中 $n = |\mathcal{T}_m|$ 为可映射任务数，$p_i \in \{0,\ldots,15\}$ 为第 $i$ 个任务（按拓扑序排列）的 PE 分配。初始种群大小 $P = 50$，采用混合策略生成：50\% 为纯随机分配，50\% 为对已有个体进行扰动（每个任务以 20\% 概率独立重新分配），确保初始多样性。

\subsection{适应度评估}

适应度评估是 GA 的计算瓶颈。每个个体需要运行一次完整的混合光电 NoC 执行所映射工作负载的 OMNeT++ 仿真。具体流程：

\begin{enumerate}
\item 生成临时 CSV 文件，写入任务到 PE 的分配；
\item 创建最小 OMNeT++ INI 文件，继承基类 \texttt{[ONoCGeneral]} 配置并指向临时 CSV；
\item 调用 \texttt{libhnocs\_dbg.exe} 子进程（Cmdenv 模式），运行完整仿真，包括 TaskPE 执行、电控制面路由、光电路建立（SETUP/ACK 握手、波长分配、MRR 调谐、SOA 放大）、数据面旁路和热网络更新；
\item 解析输出标量文件（\texttt{.sca}）：提取 makespan、每 PE DVFS 比、PE 总能耗、SOA 能耗、调谐能耗和激光器能耗；
\item 通过 \texttt{opp\_scavetool} 将输出矢量文件（\texttt{.vec}）导出为 JSON，解析每 PE 温度迹线，得到真实峰值温度 $T_{\max}$、温度标准差 $\sigma_T$ 和过热 PE 数 $N_{\text{hot}}$；
\item 基于提取的标量计算式~\eqref{eq:cost} 中的代价函数。
\end{enumerate}

超过 \SI{60}{\second} 超时阈值的仿真（表示产生了导致光层事件爆炸的病态映射）被赋予无穷大适应度，GA 自动避开搜索空间中的此类区域。

为加速评估，种群使用 \texttt{ProcessPoolExecutor} 以 8 个 worker 并行评估，每代墙钟时间减少约 6--7 倍。

\subsection{遗传算子}

\textbf{选择。}锦标赛选择，$k = 3$：从种群中不放回随机抽取三个个体，选出适应度最低者。中等选择压力下保持多样性。

\textbf{交叉。}均匀交叉，概率 $p_c = 0.8$：对每个任务，子代以等概率继承父代之一的 PE 分配。以 $1 - p_c$ 的概率直接克隆较优父代。

\textbf{变异。}逐任务变异，概率 $p_m = 0.1$：每个任务分配独立以 $p_m$ 的概率替换为随机 PE。该概率平衡了探索与利用——足够高以跳出局部最优，足够低以保留优良基因组合。

\textbf{精英保留。}$e = 2$ 个最优个体原封不动进入下一代，确保已找到的最优解永不丢失。

\textbf{终止条件。}GA 最多运行 $G = 30$ 代，若最优适应度连续 10 代无改善则提前终止。

\subsection{OMNeT++ 在环评估的必要性}

与使用简化分析热模型进行适应度评估的现有工作不同~\cite{chittamuru2015}，本方法在适应度评估和最终验证中使用\textbf{相同的} OMNeT++ 仿真器。这消除了"仿真器到仿真器"的鸿沟——即针对简化模型优化的映射在部署到全系统时性能下降的问题。OMNeT++ 模型包含：

\begin{itemize}
\item 完整光层：8$\lambda$ WDM、逐波长最低空闲优先分配、每路由器独立温度感知的 MRR 动态调谐、含饱和效应与 ASE 噪声的 SOA 放大、以及器件级光链路预算（波导交叉、弯曲、调制器、光电探测器）；
\item 热-光耦合：MRR 调谐功率注入路由器热节点，温度变化在下一时间步反馈回调谐功率；
\item 真实 DVFS：每 \SI{100}{\nano\second} 周期性重评估，含指数泄漏补偿。
\end{itemize}

\section{实验评估}

\subsection{实验设置}

在四个覆盖不同计算通信比（CCR）的基准任务图上评估所提 GA 方法：

\begin{table}[h]
\centering
\caption{基准任务图特征}
\label{tab:benchmarks}
\begin{tabular}{lcccc}
\toprule
\textbf{基准} & \textbf{任务数} & \textbf{CCR} & \textbf{结构} \\
\midrule
VOPD   & 13 & 0.3  & 长流水线 \\
MPEG4  & 12 & $\approx$1 & Fork-join + 分支 \\
HNN    & 33 & 3    & 4 级流水线 \\
GEMM   & 11 & 8    & Fork-join 归约树 \\
\bottomrule
\end{tabular}
\end{table}

GA 参数：种群大小 50，最大 30 代，锦标赛规模 3，交叉率 0.8，变异率 0.1，精英数 2，早停 patience 为 10 代，随机种子 42。所有实验在 OMNeT++ 6.3.0 + Clang 编译器 + Windows 环境下运行。基线为任务图 CSV 文件中预设的\textbf{静态映射}（由工作负载设计者启发式分配）。每个基准先运行一次基线仿真以构造式~\eqref{eq:cost} 的归一化参照，再运行 GA 搜索优化映射。

\subsection{评价指标}

报告覆盖热--性能--能耗因果链的八项指标：

\begin{itemize}
\item \textbf{热}: $T_{\max}$（\si{\degreeCelsius}，芯片峰值温度）、$\sigma_T$（K，全 PE 全时刻温度标准差）、$N_{\text{hot}}$（超过 \SI{54}{\degreeCelsius} 的 PE 数量）；
\item \textbf{性能}: $t_{\text{makespan}}$（\si{\micro\second}，任务完成时间）、$\eta_{\text{dvfs}}$（\%，平均 DVFS 节流比）；
\item \textbf{能耗}: $E_{\text{SOA}}$（\si{\micro\joule}，SOA 泵浦能耗）、$E_{\text{tune}}$（\si{\nano\joule}，MRR 动态调谐能耗）、$E_{\text{total}}$（\si{\milli\joule}，系统总能耗 = PE + SOA + 调谐 + 激光器）。
\end{itemize}

所有指标均从 OMNeT++ \texttt{.sca} 标量和 \texttt{.vec} 矢量文件中提取（使用 \texttt{opp\_scavetool} 导出），确保优化阶段与验证阶段数据一致性。

\subsection{GEMM 实验结果}

表~\ref{tab:gemm} 展示了 GEMM（CCR = 8，Fork-join 矩阵乘法归约树）的详细对比结果。GA 运行 30 代后未触发早停，耗时 37.2 分钟；因此该结果应视为当前搜索预算下的最佳发现，而非严格收敛点。

\begin{table}[h]
\centering
\caption{GEMM: 基线 vs GA 优化映射}
\label{tab:gemm}
\begin{tabular}{lccc}
\toprule
\textbf{指标} & \textbf{基线} & \textbf{GA} & \textbf{$\Delta$} \\
\midrule
$T_{\max}$ (\si{\degreeCelsius}) & 54.9 & 53.1 & $-$1.79 \\
$\sigma_T$ (K) & 2.55 & 1.89 & $-$26.0\% \\
$N_{\text{hot}}$ & 6 & 0 & $-$6 \\
$t_{\text{makespan}}$ (\si{\micro\second}) & 119.6 & 115.9 & $-$3.1\% \\
$\eta_{\text{dvfs}}$ (\%) & 1.77 & 0.00 & $-$100\% \\
$E_{\text{total}}$ (\si{\milli\joule}) & 1.569 & 1.515 & $-$3.5\% \\
$E_{\text{SOA}}$ (\si{\micro\joule}) & 1.288 & 0.541 & $-$58.0\% \\
$E_{\text{tune}}$ (\si{\nano\joule}) & 111.0 & 42.9 & $-$61.4\% \\
Comm. cost (hops$\times$bytes) & 104,448 & 62,464 & $-$40.2\% \\
Composite cost & 6.000 & 3.978 & $-$33.7\% \\
\bottomrule
\end{tabular}
\end{table}

GA 优化映射在所有指标上实现了同步改善：

\begin{itemize}
\item \textbf{热}: $T_{\max}$ 降低 \SI{1.79}{\degreeCelsius}，$\sigma_T$ 改善 26.0\%，6 个过热 PE 全部消除（$N_{\text{hot}} = 0$）。DVFS 节流比归零，说明该映射把工作负载从节流边界拉回安全区。
\item \textbf{性能}: Makespan 改善 3.1\%。由于 GEMM 的基线已经具有较高并行度，性能收益主要来自热点消除后的节流恢复，而不是任务重排带来的大幅调度改变。
\item \textbf{能耗}: 系统总能耗降低 3.5\%。光层改善更明显：SOA 能耗下降 58.0\%，调谐能耗下降 61.4\%，与通信代价下降 40.2\% 一致。
\end{itemize}

这些结果证实，在计算密集且具有归约结构的工作负载中，GA 能够同时实现通信局部化、热点削减和轻微性能提升，是本文最干净的正例。

\subsection{收敛行为}

图~\ref{fig:convergence} 展示了 GEMM 在 30 代中的适应度变化。最优适应度从第 1 代的 4.333 降至第 30 代的 3.978，改善 8.2\%；相对于基线复合代价 6.000，最终改善 33.7\%。第 27 代仍出现新的最优个体，说明 30 代预算可能仍截断了后续改进空间。

\begin{figure}[t]
\centering
% [待补充：收敛曲线图]
\caption{GEMM 30 代适应度变化。最优适应度从第 1 代 4.333 降至第 30 代 3.978；相对于基线 6.000，最终复合代价降低 33.7\%。}
\label{fig:convergence}
\end{figure}

\subsection{全基准结果总表}

表~\ref{tab:all} 汇总了全部四个基准上 GA 优化映射与静态基线的对比。

\begin{table*}[t]
\centering
\caption{四基准核心指标对比：Baseline $\rightarrow$ GA 优化映射。加粗表示改善。}
\label{tab:all}
\begin{tabular}{lccccccc}
\toprule
\textbf{基准} & \textbf{Cost} & $T_{\max}$ & $\sigma_T$ & $N_{\text{hot}}$ & Makespan & Comm. & $E_{\text{total}}$ \\
 & & (\si{\degreeCelsius}) & (K) & & (\si{\micro\second}) & (hop$\times$B) & (\si{\milli\joule}) \\
\midrule
GEMM   & 6.000$\rightarrow$\textbf{3.978} & 54.9$\rightarrow$\textbf{53.1} & 2.55$\rightarrow$\textbf{1.89} & 6$\rightarrow$\textbf{0} & 119.6$\rightarrow$\textbf{115.9} & 104448$\rightarrow$\textbf{62464} & 1.569$\rightarrow$\textbf{1.515} \\
(CCR=8) & ($-$33.7\%) & ($-$1.79\si{\degreeCelsius}) & ($-$26.0\%) & & ($-$3.1\%) & ($-$40.2\%) & ($-$3.5\%) \\
\midrule
MPEG4  & 6.000$\rightarrow$\textbf{4.035} & 54.4$\rightarrow$\textbf{52.7} & 1.54$\rightarrow$\textbf{1.46} & 2$\rightarrow$\textbf{0} & 121.7$\rightarrow$\textbf{90.9} & 420000$\rightarrow$\textbf{196000} & 1.133$\rightarrow$\textbf{0.969} \\
(CCR$\approx$1) & ($-$32.8\%) & ($-$1.69\si{\degreeCelsius}) & ($-$5.3\%) & & ($-$25.3\%) & ($-$53.3\%) & ($-$14.5\%) \\
\midrule
VOPD   & 5.000$\rightarrow$\textbf{4.063} & 52.2$\rightarrow$52.3 & 1.15$\rightarrow$1.40 & 0$\rightarrow$0 & 87.4$\rightarrow$\textbf{43.3} & 1396000$\rightarrow$\textbf{646000} & 0.747$\rightarrow$\textbf{0.510} \\
(CCR=0.3) & ($-$18.7\%) & (+0.08\si{\degreeCelsius}) & (+21.7\%) & & ($-$50.5\%) & ($-$53.7\%) & ($-$31.8\%) \\
\midrule
HNN    & 6.000$\rightarrow$\textbf{5.252} & 55.7$\rightarrow$56.0 & 3.05$\rightarrow$\textbf{2.19} & 16$\rightarrow$\textbf{8} & 204.2$\rightarrow$243.2 & 2195456$\rightarrow$\textbf{1638400} & 4.661$\rightarrow$\textbf{4.517} \\
(CCR=3) & ($-$12.5\%) & (+0.34\si{\degreeCelsius}) & ($-$28.4\%) & & (+19.1\%) & ($-$25.4\%) & ($-$3.1\%) \\
\midrule
\bottomrule
\end{tabular}
\end{table*}

\subsection{逐基准分析}

\textbf{GEMM (CCR=8, Fork-join 归约树):} 全维度正向，是四个基准中最干净的结果。GA 将通信代价削减 40.2\%，同时降低峰值温度 1.79\si{\degreeCelsius}、消除 6 个热点 PE，并使 makespan 和总能耗分别下降 3.1\% 与 3.5\%。该结果说明，对于通信结构清晰且基线存在热点的计算密集负载，任务重映射可以同时压低热点和光层代价。

\textbf{MPEG4 (CCR$\approx$1, Fork-join+分支):} 新目标函数下由旧实验中的退化案例转为强正结果。通信代价下降 53.3\%，makespan 下降 25.3\%，总能耗下降 14.5\%，且 $T_{\max}$ 下降 1.69\si{\degreeCelsius}、热点数从 2 降至 0。其 $\sigma_T$ 仅下降 5.3\%，说明收益主要来自通信路径缩短和拥塞降低，而不是大幅重塑热分布。

\textbf{VOPD (CCR=0.3, 长流水线):} 性能和能耗收益最大：makespan 下降 50.5\%，加速比约 2.02$\times$，总能耗下降 31.8\%，通信代价下降 53.7\%。但 $T_{\max}$ 略升 0.08\si{\degreeCelsius}，$\sigma_T$ 上升 21.7\%。这不是热优化失败，而是多目标优化选择了更短执行时间和更低总能耗，代价是活跃 PE 与空闲 PE 的温差变大。由于该基准全程无热点 PE，峰值温度仍低于 DVFS 阈值。

\textbf{HNN (CCR=3, 4 级流水线, 33 任务):} 这是最困难的折中案例。复合代价降低 12.5\%，通信代价降低 25.4\%，总能耗降低 3.1\%，$\sigma_T$ 改善 28.4\%，热点 PE 从 16 降至 8，平均 DVFS 惩罚从 11.01\% 降至 1.28\%。但 $T_{\max}$ 仍略升 0.34\si{\degreeCelsius}，makespan 增加 19.1\%。这说明显式 makespan 项虽然显著缓解了旧目标函数中的严重性能退化，但 30 代预算和 33 任务搜索空间仍不足以得到完全均衡的映射。

\subsection{$\sigma_T$ 的三重模式}

温度标准差 $\sigma_T$ 在四个基准上呈现三种模式，揭示了任务映射优化中热均匀性与性能、能耗之间的张力：

\begin{itemize}
\item \textbf{热-性能双赢（GEMM, MPEG4）}: $\sigma_T$ 分别下降 26.0\% 和 5.3\%，makespan 同时下降 3.1\% 和 25.3\%。通信局部化与热点削减方向一致。
\item \textbf{性能/能耗优先（VOPD）}: $\sigma_T$ 上升 21.7\%，但 makespan 和总能耗分别下降 50.5\% 与 31.8\%。在无热点约束下，提高热不均匀性换取执行时间缩短是可接受折中。
\item \textbf{未完全解决的折中（HNN）}: $\sigma_T$ 下降 28.4\%，热点 PE 减半，但 makespan 增加 19.1\%，峰值温度略升。该结果提示复杂流水线负载需要更长搜索预算或更强的并行度约束。
\end{itemize}

这一发现对 ONoC 热管理研究有直接启示：$\sigma_T$ 并非越低越好，也不能孤立解释。只有同时观察 $T_{\max}$、$N_{\text{hot}}$、makespan 和能耗，才能判断某个映射是否真正优于基线。

\subsection{收敛行为}

表~\ref{tab:convergence} 汇总了各基准的收敛统计。MPEG4 触发早停，其余均跑满 30 代。GEMM、VOPD 和 HNN 在后期仍出现新的最优个体，说明当前代数上限可能截断了进一步改进。部分个体因光层事件爆炸或仿真超时返回无穷大适应度，因此若平均适应度包含 \texttt{inf}，它不能作为种群整体质量的稳定统计；本文主要解释最终最优映射和最优适应度轨迹。

\begin{table}[h]
\centering
\caption{收敛统计}
\label{tab:convergence}
\footnotesize
\begin{tabular}{lccccc}
\toprule
\textbf{基准} & \textbf{代数} & \textbf{收敛} & \textbf{初始} & \textbf{最终} & \textbf{min} \\
\midrule
GEMM  & 30 & 否 & 4.33 & 3.98 ($-$8\%) & 37 \\
MPEG4 & 23 & 是 & 4.69 & 4.03 ($-$14\%) & 32 \\
VOPD  & 30 & 否 & 4.99 & 4.06 ($-$19\%) & 36 \\
HNN   & 30 & 否 & 6.15 & 5.25 ($-$15\%) & 63 \\
\bottomrule
\end{tabular}
\end{table}

\subsection{光层能耗分解}

光层三组分的变化模式揭示了 GA 优化作用的物理机制：

\begin{itemize}
\item \textbf{SOA 能耗}：受通信距离和电路持续时间的双重影响。GEMM 和 MPEG4 的 SOA 能耗分别下降 58.0\% 和 43.5\%，与通信代价下降一致；VOPD 与 HNN 的 SOA 能耗分别上升 13.0\% 和 5.6\%，说明缩短总通信距离并不必然降低 SOA 泵浦能耗，光路持续时间和路径重叠同样重要。
\item \textbf{调谐能耗}：GEMM、MPEG4、HNN 分别下降 61.4\%、47.7\%、36.0\%，表明减少活跃 MRR 数量和热点程度仍是降低调谐能耗的主要杠杆。VOPD 调谐能耗小幅上升 4.8\%，与其 $\sigma_T$ 变差一致。
\item \textbf{激光器能耗}：与仿真时间近似成正比。VOPD 的 makespan 下降 50.5\%，因此总能耗显著下降；HNN 的 makespan 增加 19.1\%，抵消了一部分通信和调谐收益。
\end{itemize}

\section{讨论}

\subsection{OMNeT++ 在环评估的实证价值}

本方法将全系统 OMNeT++ 仿真作为 GA 适应度评估器，其价值在 VOPD 和 HNN 两个基准上尤其明显。VOPD 的通信代价下降 53.7\%，makespan 下降 50.5\%，但 SOA 和调谐能耗并未同步下降，说明光层能耗不仅由总 hop 数决定，还受电路持续时间、路径重叠和波长槽占用影响。HNN 则相反：通信、$\sigma_T$、热点数和总能耗均改善，但 makespan 与峰值温度仍退化，暴露了复杂流水线负载中并行度保持的困难。上述现象只有在包含光路建立、波长分配、MRR 调谐、SOA 泵浦和热反馈的时域仿真中才能观察到。代价是计算开销：四个基准合计耗时约 2.8 小时，但个体间天然并行，适合设计时离线优化。

\subsection{代价函数的效果与不足}

相对于只包含峰值温度、通信、DVFS 和负载均衡的早期目标函数，本文采用的 baseline-normalized 多目标函数显著改善了 MPEG4 和 VOPD 的性能/能耗表现，并避免了 HNN 中过度牺牲 makespan 的极端结果。然而，实验也说明该函数仍存在三个不足：

\begin{enumerate}
\item \textbf{静态拥塞代理仍然粗糙}：$f_C$ 基于 XY 路径累计字节量，能够抑制空间集中通信，但不能表示同一时间窗口内的真实并发电路数。VOPD 的 SOA 能耗上升说明静态代理仍不能完全预测光层动态行为。
\item \textbf{权重对工作负载敏感}：GEMM 和 MPEG4 获得全维度改善，但 HNN 仍出现 makespan +19.1\% 和峰值温度 +0.34\si{\degreeCelsius}。对 33 任务流水线负载，$w_M$、$w_L$ 或并行度约束可能仍需提高。
\item \textbf{搜索预算限制结论强度}：GEMM、VOPD 和 HNN 在 30 代内未收敛，且后期仍出现新最优解。因此本文结果应表述为给定预算下的映射优化收益，而非全局最优。
\end{enumerate}

\subsection{$\sigma_T$ vs Makespan：热管理中的根本张力}

四个基准的结果表明，$\sigma_T$ 与 Makespan 之间不存在单调关系。GEMM 与 MPEG4 中二者同步改善；VOPD 中 $\sigma_T$ 变差但 makespan 和能耗大幅改善；HNN 中 $\sigma_T$ 改善但 makespan 退化。由此可见，$\sigma_T$ 更适合作为热均匀性约束或多目标项，而不是独立优化目标。实际系统中更合理的策略是先保证 $T_{\max}$ 和 $N_{\text{hot}}$ 不越过安全边界，再根据工作负载特性在 makespan、能耗和 $\sigma_T$ 之间选择折中点。

\subsection{局限性与未来工作}

当前研究存在若干局限：(1) 每个基准仅报告 seed=42 的一次 GA 运行，后续需要多随机种子统计显著性；(2) GEMM、VOPD 和 HNN 未在 30 代内收敛，后续应增加代数或采用代理模型预筛选以扩大搜索预算；(3) 当前拥塞项为静态 XY 字节负载代理，尚未直接记录运行时并发光电路和等待时间；(4) GA 产出的映射针对单一任务图优化，多工作负载联合优化是自然扩展；(5) MRR 因电荷俘获和热循环产生的长期谐振漂移未被建模，可靠性维度仍需后续纳入。

\section{结论}

本文提出了一种基于遗传算法的光电混合片上网络热感知任务映射框架，以全系统 OMNeT++ 仿真作为适应度评估神谕，并以静态基线归一化多目标函数联合优化热、性能、通信、拥塞、DVFS、负载和能耗。在覆盖 CCR 0.3--8 的四个基准上，核心发现如下：

\begin{enumerate}
\item \textbf{任务映射是 ONoC 热优化的有效设计时手段}。GA 优化映射在四个基准上均降低复合代价；GEMM 实现全核心指标改善（T$_{\max}$ -1.79\si{\degreeCelsius}，$N_{\text{hot}}$ 清零，makespan -3.1\%，总能耗 -3.5\%），MPEG4 同时获得热、性能、通信和能耗收益，VOPD 实现 2.02$\times$ 加速和 31.8\% 总能耗降低。

\item \textbf{OMNeT++ 在环评估是必要的}。VOPD 和 HNN 表明，通信跳数、SOA 能耗、调谐能耗、makespan 和热分布之间并不总是同向变化。光电路建立、波长争用、MRR 动态调谐和热-光耦合的完整效应需要在时域仿真中共同观察。

\item \textbf{热均匀性必须与性能和能耗联合解释}。GEMM/MPEG4 中 $\sigma_T$ 与 makespan 同步改善，VOPD 中 $\sigma_T$ 变差但性能和能耗显著改善，HNN 中 $\sigma_T$ 改善但 makespan 退化。单独最小化 $\sigma_T$ 并不等价于系统最优。

\item \textbf{复杂工作负载仍需要更强搜索和约束}。HNN 的复合代价下降 12.5\%，热点 PE 从 16 降至 8，但 makespan 仍增加 19.1\%。这提示后续工作需要多随机种子、更长代数、运行时拥塞特征和面向 DAG 并行度的约束。
\end{enumerate}

任务映射作为设计时优化手段，具有零硬件成本、与器件级和电路级热管理方案正交互补的优势。本文结果支持将其作为 ONoC 热-光-性能联合优化栈中的上层设计空间搜索方法。

\begin{thebibliography}{00}

\bibitem{batten2008}
C. Batten \emph{et al.},
``Building manycore processor-to-DRAM networks with monolithic silicon photonics,''
in \emph{Proc. IEEE Symp. High Perf. Interconnects (HOTI)}, 2008, pp. 21--30.

\bibitem{bergman2018}
K. Bergman and L. P. Carloni,
``Power evaluation of a photonic network-on-chip,''
\emph{IEEE Micro}, 2018.

\bibitem{padmaraju2014}
K. Padmaraju and K. Bergman,
``Resolving the thermal challenges for silicon microring resonator devices,''
\emph{Nanophotonics}, vol. 3, no. 4--5, pp. 269--281, 2014.

\bibitem{guha2012}
B. Guha, J. Cardenas, and M. Lipson,
``Athermal silicon microring resonators with titanium oxide cladding,''
\emph{Opt. Express}, vol. 21, pp. 26\,557--26\,563, 2012.

\bibitem{zhou2013}
L. Zhou \emph{et al.},
``Design and evaluation of an arbitration-free passive optical crossbar for on-chip interconnection networks,''
\emph{Appl. Phys. A}, vol. 95, pp. 1111--1118, 2013.

\bibitem{chittamuru2015}
S. V. R. Chittamuru, I. G. Thakkar, and S. Pasricha,
``LIBRA: Thermal and process-variation aware routing in photonic NoCs,''
\emph{IEEE Trans. Comput.-Aided Design Integr. Circuits Syst.}, vol. 37, no. 9, pp. 1778--1790, 2015.

\bibitem{skadron2003}
K. Skadron, M. R. Stan, K. Sankaranarayanan, W. Huang, S. Velusamy, and D. Tarjan,
``Temperature-aware microarchitecture: Modeling and implementation,''
\emph{ACM Trans. Archit. Code Optim.}, vol. 1, no. 1, pp. 94--125, 2003.

\end{thebibliography}

\end{document}
